\documentclass[11pt]{amsart}
\usepackage[margin=1.1in]{geometry}
\usepackage[T1]{fontenc}
\usepackage{lmodern}
\usepackage{amssymb,mathtools,microtype}
\usepackage{enumitem}
\usepackage{xcolor}
\usepackage[colorlinks=true,linkcolor=blue!50!black,citecolor=blue!50!black,urlcolor=blue!50!black]{hyperref}

\newtheorem{theorem}{Theorem}[section]
\newtheorem{proposition}[theorem]{Proposition}
\newtheorem{lemma}[theorem]{Lemma}
\newtheorem{corollary}[theorem]{Corollary}
\theoremstyle{definition}
\newtheorem{definition}[theorem]{Definition}
\theoremstyle{remark}
\newtheorem{remark}[theorem]{Remark}
\numberwithin{equation}{section}
\newcommand{\C}{\mathbb C}
\newcommand{\R}{\mathbb R}
\newcommand{\N}{\mathbb N}
\newcommand{\Tot}{\operatorname{Tot}}
\newcommand{\PSH}{\operatorname{PSH}}
\newcommand{\Levi}{\mathcal L}
\newcommand{\loc}{\mathrm{loc}}
\newcommand{\ddc}{dd^c}
\newcommand{\barp}{\bar\partial}
\newcommand{\Reg}{\mathcal R}

\title[The real analytic radius criterion]{The Radius Criterion in Real Analytic
and Lower Regularity:\\ An Expanded Proof and Audit}
\date{November 2, 2026}

\begin{document}
\begin{abstract}
For a holomorphic line bundle $A$ over a compact complex manifold,
the total space of $A^{-1}$ admits a real analytic plurisubharmonic
exhaustion if and only if $A$ admits a real analytic Hermitian metric
with semipositive curvature. We prove the equivalence, together with
its continuous, differentiable, and local H\"older versions. The proof
includes the uniform analyticity estimate needed for circle averaging,
a derivative-free positivity argument for the logarithmic radius,
and a separate verification of every regularity class. We also correct
one inference about regularity inclusions in the preceding discussion.
\end{abstract}
\maketitle
\tableofcontents

\section{Statement, conventions, and scope}

The answer to the real analytic question is \emph{yes}. The relevant
equivalence is at the \emph{same} regularity on both sides; it does not
assert that a smooth exhaustion or metric can be replaced by a real
analytic one. The following statement extends Proposition~3.3 and
Corollary~3.4 of the manuscript \cite{CX} and incorporates the
low-regularity argument discussed in the chat entitled
\emph{``Generalizing Corollary 3.4''} (the title translated from Chinese).
% Original chat title: 推广推论3.4

\begin{theorem}[Radius criterion at each regularity]\label{thm:main}
Let $Z$ be a compact complex manifold, let $A\to Z$ be a holomorphic
line bundle, and put $M=\Tot(A^{-1})$. Choose any one of the classes
\begin{equation}\label{eq:classes}
\Reg\in\bigl\{C^r:r\in\mathbb Z_{\ge0}\cup\{\infty\}\bigr\}
\ \cup\ \bigl\{C^{k,\alpha}_{\loc}:k\in\mathbb Z_{\ge0},\ 0<\alpha\le1\bigr\}
\ \cup\ \{C^\omega\}.
\end{equation}
Then the following are equivalent:
\begin{enumerate}[label=\textup{(\roman*)},leftmargin=*]
\item $M$ admits a real-valued plurisubharmonic exhaustion of class $\Reg$.
\item $A$ admits a positive definite Hermitian metric of class $\Reg$
whose local weights are plurisubharmonic.
\end{enumerate}
In particular,
\begin{equation}\label{eq:analytic-equivalence}
\begin{gathered}
\Tot(A^{-1})\text{ admits a real analytic psh exhaustion}
\\\Longleftrightarrow\\
A\text{ admits a real analytic semipositive Hermitian metric}.
\end{gathered}
\end{equation}
No K\"ahler or projectivity hypothesis on $Z$ is needed.
\end{theorem}

An \emph{exhaustion} here is a continuous real-valued function $F$
such that $\{F\le a\}$ is compact for every $a\in\R$. We use
\[
d^c=\frac{i}{4\pi}(\barp-\partial),\qquad
\ddc=\frac{i}{2\pi}\partial\barp.
\]
For a holomorphic frame $e_i$ of $A$ over $B_i\subset Z$, write
\begin{equation}\label{eq:weights}
\|e_i\|_h^2=h_i=e^{-\psi_i}>0.
\end{equation}
The metric is \emph{semipositive} if $\psi_i\in\PSH(B_i)$.
Below $C^2$ regularity this means $\ddc\psi_i\ge0$ as a current;
it does not require pointwise second derivatives. The positivity of
$h_i$ in \eqref{eq:weights} refers to the fiber norm, whereas
semipositivity refers to its curvature.

All regularity is local in real coordinates underlying holomorphic
charts. For $0<\alpha\le1$, $f\in C^{k,\alpha}_{\loc}(\Omega)$ means
$f\in C^k(\Omega)$ and, for each compact $K\Subset\Omega$,
\[
|D^\beta f(x)-D^\beta f(y)|\le C_K|x-y|^\alpha
\quad(x,y\in K,\ |\beta|=k).
\]
In particular, $C^{0,1}_{\loc}$ means locally Lipschitz. At integer
orders $C^r$ has its classical meaning; $C^{k,1}_{\loc}$ is not
identified with $C^{k+1}$. If a real noninteger index is desired, our
convention is $C^{k+\alpha}=C^{k,\alpha}_{\loc}$ for $0<\alpha<1$.
Real analyticity, denoted $C^\omega$, means representability by a
convergent power series in the real coordinates near every point.
No single analytic radius or H\"older constant on all of $M$ is required.

The metric coefficients $h_i$, their reciprocals $h_i^{-1}$, and the
weights $-\log h_i$ have the same membership in each of the classes in
\eqref{eq:classes}. Indeed, on every sufficiently small relatively
compact chart, $h_i$ is bounded away from zero, and composition with
$\log$, $\exp$, or $s\mapsto1/s$ preserves the indicated regularity.
For the analytic class these functions are themselves real analytic
on their domains.

\section{Circle averaging, including real analyticity}

Let $R_\theta(v)=e^{i\theta}v$ be the fiberwise circle action and let
$F_0$ be an exhaustion as in Theorem~\ref{thm:main}(i). Define
\begin{equation}\label{eq:average}
F(v)=\frac1{2\pi}\int_0^{2\pi}F_0(R_\theta v)\,d\theta.
\end{equation}

\begin{lemma}\label{lem:average}
The function $F$ is a circle-invariant psh exhaustion. It belongs
to every class in \eqref{eq:classes} to which $F_0$ belongs.
\end{lemma}

\begin{proof}[Proof of continuity, invariance, positivity, and properness]
For a compact set $K\subset M$, the orbit set $S^1\cdot K$ is compact.
The continuous integrand in \eqref{eq:average} is uniformly continuous
on compact coordinate neighborhoods times $S^1$; hence its integral
is continuous. Changing the variable $\theta$ shows
$F(R_\eta v)=F(v)$.

Every $R_\theta$ is holomorphic. For a holomorphic disc $f$ defined
on a neighborhood of $\overline\Delta$ in $M$, the submean inequality
for $F_0\circ R_\theta\circ f$, integrated in $\theta$, gives
\[
F(f(0))\le\frac1{2\pi}\int_0^{2\pi}F(f(e^{is}))\,ds.
\]
The interchange of integrals is valid because the integrand is
continuous on a compact set. This proves that $F$ is psh.

Put $K_a=\{F_0\le a\}$. If $F(v)\le a$, at least one point of its
circle orbit lies in $K_a$: otherwise the continuous function
$\theta\mapsto F_0(R_\theta v)$ would have minimum strictly larger
than $a$, contrary to its average. Therefore
\begin{equation}\label{eq:saturation}
\{F\le a\}\subset S^1\cdot K_a.
\end{equation}
The right side is compact, being the continuous image of
$S^1\times K_a$. The left side is closed, so it is compact as well.
Thus $F$ is an exhaustion.
\end{proof}

\begin{proof}[Proof of finite, smooth, and H\"older regularity]
In a bundle chart use real coordinates $x=(z,\lambda)$, with complex
coordinates regarded as pairs of real coordinates. Rotation is the
linear map $T_\theta(z,\lambda)=(z,e^{i\theta}\lambda)$, whose matrix
entries and operator norm are uniformly bounded in $\theta$.
For $F_0\in C^r$ and $|\beta|\le r$,
\begin{equation}\label{eq:diff-average}
D^\beta F(x)=\frac1{2\pi}\int_0^{2\pi}
D_x^\beta(F_0\circ T_\theta)(x)\,d\theta.
\end{equation}
The derivatives of the integrand are jointly continuous and uniformly
bounded on the relevant compact sets. Difference quotients and
dominated convergence justify \eqref{eq:diff-average}. Its right side
is continuous. Applying this for every finite $r$ proves the
$C^\infty$ assertion as well.

For $F_0\in C^{k,\alpha}_{\loc}$, the chain rule for the linear maps
$T_\theta$ and the local H\"older estimates on the compact orbit set
give, on a smaller compact coordinate neighborhood,
\[
\bigl|D^\beta(F_0\circ T_\theta)(x)
-D^\beta(F_0\circ T_\theta)(y)\bigr|
\le C|x-y|^\alpha,\qquad |\beta|=k,
\]
with $C$ independent of $\theta$. A finite neighborhood cover of the
orbit set supplies such a uniform constant; pairs of points not close
enough to lie in one of these neighborhoods are handled by boundedness
of the derivatives. Integrating the estimate proves the same H\"older
bound for $D^\beta F$. This also applies when $k=0$ or $\alpha=1$.
\end{proof}

\begin{lemma}[Analytic integration over a compact parameter]\label{lem:analytic-integral}
Suppose $G(x,\theta)$ is jointly real analytic on a neighborhood of
$\{x_0\}\times S^1$, where $x\in\R^m$. Then
\[
I(x)=\frac1{2\pi}\int_0^{2\pi}G(x,\theta)\,d\theta
\]
is real analytic near $x_0$.
\end{lemma}

\begin{proof}
Use a real angular coordinate near each point of $S^1$. The convergent
real power series of $G$ in $(x,\theta)$ extends holomorphically to
the complexified variables on a smaller polydisc. After shrinking
that polydisc, its extension is bounded there. Consequently, for each
angular point there is an arc $J$, a number $R_J>0$, and a bound $M_J$
such that, for every real $\theta\in J$, the function
$x\mapsto G(x,\theta)$ has a bounded holomorphic extension to
\[
|x_j-x_{0,j}|<2R_J,\qquad 1\le j\le m.
\]
Compactness of $S^1$ supplies finitely many such arcs. Taking a smaller
common radius and the largest of the finitely many bounds, Cauchy's
estimates give constants $R,M>0$ such that
\begin{equation}\label{eq:analytic-estimate}
\left|\frac{D_x^\beta G(x_0,\theta)}{\beta!}\right|
\le M R^{-|\beta|}
\quad\text{for every multi-index $\beta$ and every $\theta$.}
\end{equation}
Each local extension has the given real function as its restriction,
so its Taylor coefficients in $x$ are precisely those in
\eqref{eq:analytic-estimate}. No choice of a global complexification
in the angular variable is needed.

For $|x-x_0|_\infty<R$, the Taylor series
\[
G(x,\theta)=\sum_{\beta\in\mathbb Z_{\ge0}^m}
\frac{D_x^\beta G(x_0,\theta)}{\beta!}(x-x_0)^\beta
\]
converges uniformly in $\theta$ on every smaller polydisc. Indeed,
if $|x-x_0|_\infty\le r<R$, the sum of the absolute bounds is at most
\[
M\sum_\beta(r/R)^{|\beta|}=M(1-r/R)^{-m}<\infty.
\]
Termwise integration is therefore justified and yields the convergent
power series
\[
I(x)=\sum_\beta\left(
\frac1{2\pi}\int_0^{2\pi}
\frac{D_x^\beta G(x_0,\theta)}{\beta!}\,d\theta\right)(x-x_0)^\beta.
\]
Thus $I$ is real analytic near $x_0$.
\end{proof}

\begin{proof}[Analytic case of Lemma~\ref{lem:average}]
In a bundle chart the function
\[
G(x,\theta)=F_0(T_\theta x)
\]
is jointly real analytic: the rotation
matrix has analytic entries $\sin\theta$ and $\cos\theta$, and
$F_0$ is real analytic. Apply Lemma~\ref{lem:analytic-integral} at
each $x_0$. This proves $F\in C^\omega$, including at the zero section.
\end{proof}

\begin{remark}
Differentiating the integral to every order proves smoothness, but
does not by itself prove analyticity. The uniform convergence radius,
or equivalently the factorial estimates \eqref{eq:analytic-estimate},
is the additional ingredient. There is also no appeal to approximation
of smooth psh functions by analytic ones.
\end{remark}

\section{Logarithmic radii from a continuous exhaustion}

Fix a holomorphic frame $e_i$ of $A$ on a coordinate neighborhood
$B_i$. Write vectors of $A^{-1}$ as $v=\lambda_i e_i^*(z)$ and set
\begin{equation}\label{eq:UV}
U_i(z,t)=F\bigl(e^{t/2}e_i^*(z)\bigr),\qquad
V_i(z,w)=F\bigl(e^{w/2}e_i^*(z)\bigr).
\end{equation}
The map $(z,w)\mapsto e^{w/2}e_i^*(z)$ is holomorphic. Consequently
$V_i$ is psh on $B_i\times\C$, and circle invariance gives
\begin{equation}\label{eq:real-part}
V_i(z,w)=U_i(z,\operatorname{Re}w).
\end{equation}

\subsection{Convexity without classical second derivatives}

For fixed $z$, write $u(t)=U_i(z,t)$ and $w=t+is$.
Subharmonicity of $u(\operatorname{Re}w)$ implies, as distributions,
\[
\Delta_w\bigl(u(t)\bigr)=u''(t)\ge0.
\]
To read this as an inequality on the real line, test against a product
$\eta(t)\chi(s)$ with $\eta,\chi\ge0$ and $\int\chi=1$; the term
involving $\chi''$ integrates to zero. Thus
$\langle u'',\eta\rangle\ge0$.
Convolve $u$ on compact subintervals with a nonnegative smooth
approximate identity. The resulting smooth functions satisfy
$u_\varepsilon''\ge0$ and are convex. They converge locally uniformly
to the continuous function $u$. Passing to the limit in the convexity
inequality proves that $u$ is convex.

Moreover,
\begin{equation}\label{eq:limits}
\lim_{t\to-\infty}U_i(z,t)=F(0_z),\qquad
\lim_{t\to+\infty}U_i(z,t)=+\infty.
\end{equation}
The first limit is continuity at the zero vector. If the second failed,
there would be a sequence $t_\nu\to+\infty$ along which $U_i(z,t_\nu)$
is bounded above. Its corresponding vectors would lie in one compact
sublevel of $F$, although their fiber coordinates tend to infinity.
This is impossible.

\begin{lemma}[Unique upper-level radius]\label{lem:convex}
Let $u:\R\to\R$ be continuous and convex, with
$\lim_{t\to-\infty}u(t)=\ell\in\R$ and
$\lim_{t\to+\infty}u(t)=+\infty$. Then $u$ is nondecreasing.
For every $c>\ell$ there is a unique $s\in\R$ such that $u(s)=c$,
and
\[
u(t)<c\ \Longleftrightarrow\ t<s,\qquad
u(t)\le c\ \Longleftrightarrow\ t\le s.
\]
If $u$ is differentiable at $s$, then $u'(s)>0$.
\end{lemma}

\begin{proof}
If $u(b)<u(a)$ for some $a<b$, let
$d=(u(b)-u(a))/(b-a)<0$. The monotonicity of secant slopes for convex
functions implies, for $x<a$,
\[
\frac{u(a)-u(x)}{a-x}\le d,
\qquad u(x)\ge u(a)-d(a-x)\longrightarrow+\infty
\quad(x\to-\infty),
\]
contradicting the finite limit. Hence $u$ is nondecreasing, and $u\ge\ell$.
The intermediate value theorem supplies a point at level $c$.
If $u(s)=u(t)=c$ with $s<t$, monotonicity gives
$u(x)\le c$ for $x<s$, while convexity gives
\[
0\le\frac{u(s)-u(x)}{s-x}\le
\frac{u(t)-u(s)}{t-s}=0.
\]
Thus $u(x)=c$ for all $x<s$, contradicting $\ell<c$.
Uniqueness together with monotonicity gives both stated equivalences.
Finally, for any $a<s$, the supporting-slope inequality gives
\begin{equation}\label{eq:positive-derivative}
u'(s)\ge\frac{u(s)-u(a)}{s-a}
=\frac{c-u(a)}{s-a}>0.
\end{equation}
\end{proof}

\subsection{Choice and continuity of the radius}

Since the zero section is compact, choose once and for all
\begin{equation}\label{eq:c}
c>\max_{z\in Z}F(0_z).
\end{equation}
Lemma~\ref{lem:convex} provides a unique function $\sigma_i:B_i\to\R$
such that
\begin{equation}\label{eq:radius-equation}
U_i(z,\sigma_i(z))=c.
\end{equation}
For nonzero vectors, $t=\log|\lambda_i|^2$, so
\begin{equation}\label{eq:disc-bundle}
\{F<c\}\cap A^{-1}|_{B_i}
=\{(z,\lambda_i):|\lambda_i|^2<e^{\sigma_i(z)}\}.
\end{equation}
The zero vectors belong to both sides by \eqref{eq:c}. Thus
$e^{\sigma_i(z)/2}$ is exactly the fiber radius of this sublevel set.

Continuity does not require an implicit function theorem. Fix $z_0$,
write $s_0=\sigma_i(z_0)$, and take $\varepsilon>0$. The strict
inequalities
\[
U_i(z_0,s_0-\varepsilon)<c<U_i(z_0,s_0+\varepsilon)
\]
persist for $z$ near $z_0$ by continuity of $U_i$. The unique level root
must then lie between $s_0-\varepsilon$ and $s_0+\varepsilon$.
Therefore $\sigma_i$ is continuous.

\section{Positivity of the negative logarithmic radius}

\begin{proposition}\label{prop:psh-radius}
For the continuous radii constructed above, $\psi_i=-\sigma_i$ is psh.
\end{proposition}

\begin{proof}
Let $f$ be a holomorphic map from a neighborhood of $\overline\Delta$
to $B_i$. The function
$\varphi(e^{i\theta})=\sigma_i(f(e^{i\theta}))$ is continuous on
$\partial\Delta$. Its Poisson integral is
\begin{equation}\label{eq:poisson}
H(re^{it})=\frac1{2\pi}\int_0^{2\pi}
\frac{1-r^2}{1-2r\cos(t-\theta)+r^2}
\varphi(e^{i\theta})\,d\theta,\qquad 0\le r<1.
\end{equation}
The Poisson kernel is nonnegative, has integral $2\pi$, and
concentrates at its boundary point as $r\uparrow1$. Uniform continuity
of $\varphi$ therefore gives a continuous extension of $H$ to
$\overline\Delta$, equal to $\varphi$ on the boundary.
The Poisson kernel is harmonic in the interior variable, so $H$ is
harmonic and smooth on $\Delta$. In particular,
\begin{equation}\label{eq:mean-H}
H(0)=\frac1{2\pi}\int_0^{2\pi}\sigma_i(f(e^{i\theta}))\,d\theta.
\end{equation}

Writing the disc coordinate as $x+iy$, the function $H_x-iH_y$ is
holomorphic because $H_{xx}+H_{yy}=0$. Since $\Delta$ is simply
connected, it has a holomorphic primitive $g$. Adjusting the real
constant of integration gives $\operatorname{Re}g=H$.
The map
\[
\zeta\longmapsto e^{g(\zeta)/2}e_i^*(f(\zeta))
\]
is holomorphic on $\Delta$. Hence
\begin{equation}\label{eq:q}
q(\zeta)=F\bigl(e^{g(\zeta)/2}e_i^*(f(\zeta))\bigr)
=U_i(f(\zeta),H(\zeta))
\end{equation}
is subharmonic there. The second expression is continuous on
$\overline\Delta$, and its boundary value is
\[
U_i\bigl(f(\zeta),\sigma_i(f(\zeta))\bigr)=c.
\]
The maximum principle yields $q\le c$. By Lemma~\ref{lem:convex},
this is equivalent to
\[
H(\zeta)\le\sigma_i(f(\zeta)).
\]
At the center, using \eqref{eq:mean-H} and multiplying by $-1$, we get
\[
-\sigma_i(f(0))\le\frac1{2\pi}\int_0^{2\pi}
-\sigma_i(f(e^{i\theta}))\,d\theta.
\]
Applying the same argument to every smaller disc proves the
submean inequalities for $-\sigma_i\circ f$. Since $-\sigma_i$ is
continuous, the holomorphic-disc criterion proves it is psh.
\end{proof}

\begin{remark}\label{rem:conjugate}
The harmonic conjugate of continuous boundary data need not extend
continuously to the boundary. The proof does not assume this:
$q=U_i(f,H)$ has the required continuous extension even when $g$
does not. Circle invariance is precisely what removes the imaginary
part of $g$ from this boundary argument.
\end{remark}

\subsection{Gluing and the sign of the curvature}

On an overlap, suppose $e_j=g_{ij}e_i$, with $g_{ij}$ nowhere-vanishing
and holomorphic. Then $e_j^*=g_{ij}^{-1}e_i^*$, and equality of the
two descriptions of a vector gives
\[
\lambda_i e_i^*=\lambda_j e_j^*,\qquad
\lambda_j=g_{ij}\lambda_i.
\]
At the level $F=c$, \eqref{eq:disc-bundle} therefore gives
\begin{equation}\label{eq:transition}
e^{\sigma_j}=|g_{ij}|^2e^{\sigma_i},\quad
\sigma_j=\sigma_i+\log|g_{ij}|^2,\quad
\psi_j=\psi_i-\log|g_{ij}|^2.
\end{equation}
Consequently
\begin{equation}\label{eq:constructed-metric}
\|e_i\|_h^2=e^{-\psi_i}=e^{\sigma_i}
\end{equation}
defines a positive definite continuous Hermitian metric on $A$.
Its curvature is $\ddc\psi_i\ge0$. The curvature currents agree
because $\log|g_{ij}|^2$ is pluriharmonic locally. This completes the
continuous forward implication, with no differentiation of $F$.

\section{Recovering precisely the original regularity}

Composition in \eqref{eq:UV} preserves every class in
\eqref{eq:classes}, since $(z,t)\mapsto e^{t/2}e_i^*(z)$ is real
analytic. We now recover the regularity of $\sigma_i$ from the
level equation. The three cases below require different justifications.

\subsection{Classical \texorpdfstring{$C^r$}{C-r} and real analytic regularity}

Suppose first that $U_i$ is $C^1$. At each level root,
Lemma~\ref{lem:convex} gives
\begin{equation}\label{eq:transversality}
\partial_tU_i(z,\sigma_i(z))>0.
\end{equation}
For any finite $r\ge1$, the $C^r$ implicit function theorem applied
to $U_i(z,t)-c=0$ therefore gives $\sigma_i\in C^r$ when $U_i\in C^r$.
The local solution agrees with the previously constructed continuous
radius by uniqueness. Applying this for every finite $r$ gives the
$C^\infty$ conclusion.

When $U_i\in C^\omega$, the \emph{real analytic} implicit function
theorem applies. More explicitly, take real base coordinates
$x\in\R^{2\dim_\C Z}$ near $z_0$ and put $t_0=\sigma_i(z_0)$.
Complexify the convergent power series of $U_i(x,t)-c$ near
$(x_0,t_0)$. The resulting holomorphic function
$\widetilde U(X,T)-c$ satisfies
$\partial_T\widetilde U(x_0,t_0)\ne0$. The holomorphic map
\[
(X,T)\longmapsto (X,\widetilde U(X,T)-c)
\]
has invertible Jacobian there. Its local holomorphic inverse, restricted
to the second coordinate $0$, supplies a holomorphic function
$T=S(X)$ solving the level equation. The complexified series is real
on real inputs; uniqueness, or conjugation of this equation, implies
that $S(x)$ is real for real $x$ near $x_0$. Thus its restriction is
real analytic and equals $\sigma_i$ by uniqueness. Formula
\eqref{eq:constructed-metric} now makes $h$ real analytic.

\begin{remark}
No regular-value selection by Sard's theorem is needed. Every
$c>\max_ZF(0_z)$ has the required positive radial derivative. In the
$C^1$ case, this also implies that $F=c$ is a regular level: along
$s\mapsto e^s v$ its derivative is $2U_{i,t}>0$ at such a point.
\end{remark}

\subsection{The \texorpdfstring{$C^{0,\alpha}$}{C-(0,alpha)} case, including Lipschitz regularity}

The implicit function theorem is unavailable at this regularity.
Fix $z_0$ and abbreviate $U=U_i$, $\sigma=\sigma_i$. Choose
\[
a<b<\sigma(z_0)<d,\qquad U(z_0,b)>U(z_0,a).
\]
Such $a,b$ exist: by continuity there is a $b<\sigma(z_0)$ with
$U(z_0,b)>F(0_{z_0})$, and then the limit at $-\infty$ supplies $a<b$
with the strict inequality. Continuity of $U$ and $\sigma$ permits
shrinking the base neighborhood so that
\begin{equation}\label{eq:secant}
b<\sigma(z)<d,\qquad
\frac{U(z,b)-U(z,a)}{b-a}\ge m>0
\end{equation}
throughout it, with one constant $m$. Ordered secant slopes of a
convex function now imply
\begin{equation}\label{eq:strong-increase}
U(z,t)-U(z,s)\ge m(t-s),\qquad b\le s\le t\le d.
\end{equation}

Assume $\sigma(z)\ge\sigma(w)$. Since both radii lie in $(b,d)$,
\begin{align*}
m|\sigma(z)-\sigma(w)|
&\le U(z,\sigma(z))-U(z,\sigma(w))\\
&=U(w,\sigma(w))-U(z,\sigma(w))\\
&\le C|z-w|^\alpha.
\end{align*}
The last estimate follows from the $C^{0,\alpha}$ bound for $U$ on
a compact product of a smaller base neighborhood with $[b,d]$.
If the two radii have the opposite order, interchange $z,w$. Hence
\begin{equation}\label{eq:holder-radius}
|\sigma(z)-\sigma(w)|\le(C/m)|z-w|^\alpha.
\end{equation}
This proves the same H\"older exponent for the radius, including
$\alpha=1$. Taking its exponential preserves that exponent for $h$.

\subsection{The \texorpdfstring{$C^{k,\alpha}$}{C-(k,alpha)} case for \texorpdfstring{$k\ge1$}{k at least 1}}

For completeness, one can obtain the H\"older version of the implicit
function theorem here directly. The classical theorem first gives
$\sigma\in C^k$. On a small convex real coordinate neighborhood,
$\sigma$ is Lipschitz because its first derivatives are bounded.
Write $x$ for the real base variables. Differentiating $U(x,\sigma(x))=c$
once yields
\begin{equation}\label{eq:first-derivative}
D_{x_j}\sigma(x)=-\frac{U_{x_j}(x,\sigma(x))}{U_t(x,\sigma(x))}.
\end{equation}
The denominator has a positive lower bound after shrinking the
neighborhood, by \eqref{eq:transversality}. For $k=1$, the first
derivatives of $U$ are $\alpha$-H\"older, and composition with the
Lipschitz map $x\mapsto(x,\sigma(x))$ preserves this property.
Products and reciprocals of bounded H\"older functions with nonzero
denominator preserve it as well. Thus $D\sigma$ is $\alpha$-H\"older.

For $k\ge2$, differentiating the level equation $k$ times gives,
for each multi-index $\beta$ of length $k$,
\begin{equation}\label{eq:higher-derivative}
U_t(x,\sigma(x))D^\beta\sigma(x)=-P_\beta(x).
\end{equation}
Here $P_\beta$ is a finite sum of products of derivatives of $U$ of
order at most $k$, evaluated at $(x,\sigma(x))$, and derivatives of
$\sigma$ of order at most $k-1$. This form follows inductively from
the chain and product rules: the only term involving a derivative of
$\sigma$ of order $k$ is the displayed $U_tD^\beta\sigma$.
Derivatives of $\sigma$ of order at most $k-1$ are locally Lipschitz
since $\sigma\in C^k$. Derivatives of $U$ of order less than $k$ are
locally Lipschitz, and those of order $k$ are $\alpha$-H\"older.
Their compositions with $(x,\sigma(x))$ therefore have the same bounds.
All factors of $P_\beta$ are bounded and $\alpha$-H\"older on the
smaller neighborhood, so $P_\beta$ is $\alpha$-H\"older. Division by
$U_t$ in \eqref{eq:higher-derivative} proves that $D^\beta\sigma$ is
$\alpha$-H\"older. This argument works unchanged at $\alpha=1$.

Together with the continuous construction, these arguments prove
Theorem~\ref{thm:main}(i)$\Rightarrow$(ii) for every listed class.

\section{From a semipositive metric to an exhaustion}

Let $h$ be a metric as in Theorem~\ref{thm:main}(ii), with local
weights $\psi_i$. Define the squared dual norm
\begin{equation}\label{eq:rho}
\rho(v)=\|v\|_{h^{-1}}^2.
\end{equation}
Since $\|e_i^*\|_{h^{-1}}^2=1/\|e_i\|_h^2=e^{\psi_i}$, its local
formula is
\begin{equation}\label{eq:rho-local}
\rho(z,\lambda_i)=|\lambda_i|^2e^{\psi_i(z)}.
\end{equation}
Under the frame changes \eqref{eq:transition},
\[
|\lambda_j|^2e^{\psi_j}
=|g_{ij}\lambda_i|^2e^{\psi_i-\log|g_{ij}|^2}
=|\lambda_i|^2e^{\psi_i}.
\]
Thus these formulas agree on overlaps.

\subsection{Regularity, including the zero section}

The function $|\lambda_i|^2$ is a polynomial in the real fiber
coordinates. Composition of $\psi_i$ with the exponential and
multiplication by this polynomial preserve $C^r$, $C^{k,\alpha}_{\loc}$,
and $C^\omega$. Therefore $\rho$ has the stated regularity on the
entire bundle chart, including $\lambda_i=0$. There is no division
by $\lambda_i$ or logarithmic singularity in \eqref{eq:rho-local}.
For example, for bounded H\"older functions the elementary estimate
\[
|a(x)b(x)-a(y)b(y)|
\le\|a\|_\infty|b(x)-b(y)|+\|b\|_\infty|a(x)-a(y)|
\]
gives the product bound; applying the product and chain rules gives
the higher-order versions. On this noncompact total space, the
result is local regularity, not a uniform global H\"older bound.

\subsection{Plurisubharmonicity with no loss at zero}

Away from the zero section,
\[
\log\rho=\log|\lambda_i|^2+\psi_i(z)
\]
is psh, because the first summand is pluriharmonic there. The
exponential is increasing and convex, so $\rho$ is psh there.
Indeed, for a subharmonic function $u$, its submean inequality and
Jensen's inequality give
\[
e^{u(0)}\le e^{\frac1{2\pi}\int u(e^{it})dt}
\le\frac1{2\pi}\int e^{u(e^{it})}\,dt.
\]
The same argument on holomorphic discs proves the psh assertion.

Here is an explicit local proof across the zero section, avoiding
any regularity assumption on second derivatives of $\psi_i$.
For $\varepsilon>0$ set, on a fixed bundle chart,
\[
\rho_\varepsilon(z,\lambda_i)
=(|\lambda_i|^2+\varepsilon)e^{\psi_i(z)}.
\]
The function
$\log\rho_\varepsilon=\log(|\lambda_i|^2+\varepsilon)+\psi_i(z)$
is psh, since
\[
\frac{\partial^2}{\partial\lambda_i\partial\bar\lambda_i}
\log(|\lambda_i|^2+\varepsilon)
=\frac{\varepsilon}{(|\lambda_i|^2+\varepsilon)^2}\ge0.
\]
Thus $\rho_\varepsilon$ is psh by convex composition.
Moreover $\rho_\varepsilon\to\rho$ locally uniformly, since
$\rho_\varepsilon-\rho=\varepsilon e^{\psi_i}$ and $\psi_i$ is
continuous. Passing to the limit in the disc submean inequalities
shows that $\rho$ is psh throughout the chart. These auxiliary
$\rho_\varepsilon$ need not agree on overlaps: they are used only
to prove the local property of the globally defined $\rho$.

Equivalently, one may use the removable-singularity theorem for
locally bounded psh functions across the zero section. The local
argument above verifies directly that the extension has value zero.

\subsection{Compactness of the sublevels}

Choose any smooth reference Hermitian metric $k$ on $A^{-1}$.
Since the fibers are one-dimensional, the ratio
\[
a(z)=\frac{\|v\|_{h^{-1}}^2}{\|v\|_k^2},\qquad 0\ne v\in A^{-1}_z,
\]
is independent of the choice of $v$. It is a positive continuous
function on compact $Z$. Hence there are constants $0<m\le M<\infty$
such that
\begin{equation}\label{eq:comparison}
m\|v\|_k^2\le\rho(v)\le M\|v\|_k^2.
\end{equation}
The closed disk bundle $\{\|v\|_k^2\le R\}$ over a compact base is
compact. To check this explicitly, cover $Z$ by finitely many compact
sets $K_j$ each contained in a bundle chart. On $K_j$, the positive
coefficient of $k$ has a positive minimum, so the indicated disk
bundle lies in $K_j\times\{|\lambda|\le C_j\}$ and is closed there.
Each such piece is compact, and their finite union is compact.

For $b\ge0$, \eqref{eq:comparison} places the closed set
$\{\rho\le b\}$ inside the compact disk bundle
$\{\|v\|_k^2\le b/m\}$. For $b<0$ the sublevel is empty.
Thus $\rho$ is an exhaustion. This proves
Theorem~\ref{thm:main}(ii)$\Rightarrow$(i), and completes the theorem.

\begin{remark}[Geometric normalization]
If $h$ was obtained from a circle-averaged exhaustion and a level $c$,
then $\psi_i=-\sigma_i$, so
\[
\rho(z,\lambda_i)=|\lambda_i|^2e^{-\sigma_i(z)},\qquad
\{\rho<1\}=\{F<c\}.
\]
The new exhaustion need not equal $F$; it is the squared Hermitian
gauge of this one complete circular sublevel set.
\end{remark}

\section{The classical Levi calculations, written out}

The preceding proof establishes positivity at continuous regularity.
For comparison with the original Proposition~3.3, we verify its
$C^2$ calculation, including the sign and the factor $2$.

\subsection{Levi form of the radius}

Assume $U$ and $\sigma$ are $C^2$ and omit the chart index. All
derivatives of $U$ below are evaluated at $(z,\sigma(z))$.
Differentiating $U(z,\sigma(z))=c$ gives
\[
U_{z_j}+U_t\sigma_{z_j}=0
\]
and then
\begin{equation}\label{eq:second-radius}
0=U_{z_j\bar z_k}+U_{z_jt}\sigma_{\bar z_k}
+U_{t\bar z_k}\sigma_{z_j}
+U_{tt}\sigma_{z_j}\sigma_{\bar z_k}
+U_t\sigma_{z_j\bar z_k}.
\end{equation}
For $\xi\in T_z^{1,0}Z$, put $b=\partial\sigma(\xi)$.
Since $V(z,w)=U(z,\operatorname{Re}w)$,
\[
V_{z_j\bar w}=\tfrac12U_{z_jt},\qquad
V_{w\bar z_k}=\tfrac12U_{t\bar z_k},\qquad
V_{w\bar w}=\tfrac14U_{tt}.
\]
At the point $(z,w=\sigma(z))$ this gives
\begin{align*}
\Levi_V(\xi,2b)
&=\sum_{j,k}U_{z_j\bar z_k}\xi_j\bar\xi_k
+2\operatorname{Re}\left(\sum_jU_{z_jt}\xi_j\bar b\right)
+U_{tt}|b|^2\\
&=-U_t\Levi_\sigma(\xi),
\end{align*}
where the last equality is \eqref{eq:second-radius} contracted with
$\xi_j\bar\xi_k$. Because $V$ is psh and $U_t>0$,
\begin{equation}\label{eq:Levi-radius}
\Levi_{-\sigma}(\xi)=\frac{\Levi_V(\xi,2\partial\sigma(\xi))}{U_t}
\ge0.
\end{equation}
The plus sign in the vector $(\xi,2\partial\sigma(\xi))$ is correct.

\subsection{Levi form of the squared dual norm}

For $\rho=|\lambda|^2e^\psi$, direct differentiation yields
\begin{align*}
\rho_{z_j\bar z_k}
&=|\lambda|^2e^\psi(\psi_{z_j\bar z_k}
+\psi_{z_j}\psi_{\bar z_k}),\\
\rho_{z_j\bar\lambda}&=\lambda e^\psi\psi_{z_j},\qquad
\rho_{\lambda\bar z_k}=\bar\lambda e^\psi\psi_{\bar z_k},\qquad
\rho_{\lambda\bar\lambda}=e^\psi.
\end{align*}
Contracting with $(\xi,\eta)$ and completing the square gives
\begin{equation}\label{eq:Levi-norm}
\Levi_\rho(\xi,\eta)
=e^\psi\left(
|\eta+\lambda\partial\psi(\xi)|^2
+|\lambda|^2\Levi_\psi(\xi)\right)\ge0.
\end{equation}
At $\lambda=0$ this reduces to $e^\psi|\eta|^2$, so it also explicitly
checks positivity on the zero section in the analytic and $C^2$ cases.

\section{Audit of the earlier discussion}

The continuous radius argument and its H\"older extension are valid.
The following qualifications distinguish the established conclusions
from stronger statements that do not follow from this proof.

\begin{enumerate}[label=\textup{\arabic*.},leftmargin=*]
\item \textbf{Analytic averaging.}
The compactness estimate in Lemma~\ref{lem:analytic-integral} supplies
uniform Taylor convergence. Smooth differentiation under the integral
alone is insufficient. With this addition, the analytic implicit
function theorem proves the requested real analytic equivalence.

\item \textbf{Low-regularity positivity is not a Hessian computation.}
For merely continuous or $C^1$ data, Proposition~\ref{prop:psh-radius}
proves the psh property by harmonic comparison. Curvature is interpreted
as a positive current. The second-derivative identity
\eqref{eq:Levi-radius} is used only when those derivatives exist.

\item \textbf{There is no hidden boundary assumption on a harmonic conjugate.}
Only $H$ and the expression $U_i(f,H)$ must extend continuously to the
closed disc; $g$ need not do so. See Remark~\ref{rem:conjugate}.

\item \textbf{H\"older regularity at order zero uses an inverse estimate.}
The strictly positive uniform secant slope in \eqref{eq:secant}, not
an unavailable differentiable implicit function theorem, is what
preserves $C^{0,\alpha}$ and Lipschitz regularity. For $k\ge1$,
\eqref{eq:first-derivative}--\eqref{eq:higher-derivative} verify the
H\"older conclusion, including the endpoint $\alpha=1$.

\item \textbf{An equivalence at stronger regularity does not descend by inclusion alone.}
Let $E_{\Reg}$ denote existence of an exhaustion of class $\Reg$, and
$H_{\Reg}$ existence of a metric of that class. If
$\Reg_s\subset\Reg_t$, then
\[
E_{\Reg_s}\Longrightarrow E_{\Reg_t},\qquad
H_{\Reg_s}\Longrightarrow H_{\Reg_t}.
\]
Even together with $E_{\Reg_s}\Longleftrightarrow H_{\Reg_s}$, these
arrows do not imply $E_{\Reg_t}\Longleftrightarrow H_{\Reg_t}$:
the two weaker existence properties may enlarge independently.
The earlier chat's claim that the weaker equivalence follows
automatically by inclusion must therefore be removed. The direct
proof above establishes the equivalence separately at every class.

\item \textbf{The H\"older inclusions themselves need the stated endpoint convention.}
For $0<\alpha<\beta\le1$,
$C^{k,\beta}_{\loc}\subset C^{k,\alpha}_{\loc}$, since on a compact
set of diameter $D$,
\[
|x-y|^\beta\le D^{\beta-\alpha}|x-y|^\alpha.
\]
Bounded first derivatives on smaller coordinate balls also imply
\[
C^{k+1}\subsetneq C^{k,1}_{\loc}\subset C^{k,\alpha}_{\loc}
\quad(0<\alpha<1).
\]
For the strict first inclusion, use $|x|$ when $k=0$; for $k\ge1$,
take a $k$-fold antiderivative of $|x|$, whose $k$th derivative is
Lipschitz but not differentiable at zero. The example
$|x|^{3/2}\in C^{1,1/2}\setminus C^2$ illustrates a different
endpoint issue and does not by itself prove
$C^{k+1}\subsetneq C^{k,1}_{\loc}$. For noninteger orders, the usual
ordering by $k+\alpha$ is valid; integer differentiability and
Lipschitz endpoints must remain distinct.

\item \textbf{Existence at every regularity is not asserted.}
For each class, the theorem equates two existence questions for a
fixed bundle. It does not assert that any of them has a positive
answer, nor that $C^\infty$ semipositivity implies $C^\omega$
semipositivity. Likewise, nonexistence of a $C^2$ metric rules out
$C^2$, smooth, and analytic exhaustions, but does not alone rule out
$C^1$ or locally Lipschitz exhaustions. In the particular example of
\cite[Theorem~1.1]{CX}, the $C^2$ obstruction already rules out the
analytic class; the lower-regularity threshold requires its own argument.

\item \textbf{The hypotheses and signs are essential to the construction.}
Compactness of $Z$ supplies one level above the whole zero section
and compact norm disk bundles. The object constructed is a metric
on $A$, while the exhaustion lives on $A^{-1}$. With the convention
$\|e_i\|_h^2=e^{-\psi_i}$, the psh weight is $-\sigma_i$, and the
metric coefficient is $e^{\sigma_i}$. These signs and the factor $2$
in \eqref{eq:Levi-radius} agree with the original $C^2$ proof.
\end{enumerate}

\begin{corollary}[Replacement for Corollary 3.4]\label{cor:replacement}
Let $A\to Z$ be a holomorphic line bundle over a compact complex
manifold. For each $r\in\mathbb Z_{\ge0}\cup\{\infty,\omega\}$,
$\Tot(A^{-1})$ admits a $C^r$ psh exhaustion if and only if $A$
admits a $C^r$ Hermitian metric with psh local weights. The same
equivalence holds with $C^r$ replaced on both sides by
$C^{k,\alpha}_{\loc}$, for every $k\in\mathbb Z_{\ge0}$ and $0<\alpha\le1$.
Here $C^\omega$ means real analytic, and below $C^2$ semipositivity
is understood in the sense of currents.
\end{corollary}

\paragraph{Source and scope note.}
The audit concerns the radius criterion and the regularity discussion
in the referenced chat, not the other geometric arguments in the
manuscript. The proofs of the averaging, radius, and regularity steps
are supplied in this note. Standard background on psh functions,
holomorphic pullback, convex composition, and removable singularities
can be found in \cite[Chapter~I, \S5]{Demailly}; in that chapter see
Definition~5.1 and Theorems~5.6, 5.11, and 5.24.

\begin{thebibliography}{9}
\bibitem{CX}
Z.~Chen and B.~Xiao,
\emph{Continuous and Smooth Plurisubharmonic Exhaustions:
Threefold Counterexamples and Surface Regularization},
manuscript, October 30, 2026, \S3, especially Proposition~3.3 and
Corollary~3.4. The version consulted is the local source
\texttt{10-30-revision/psh-exhaustion.tex}.

\bibitem{Demailly}
J.-P.~Demailly,
\emph{Complex Analytic and Differential Geometry},
open-content book, Chapter~I, \S5.
\href{https://www-fourier.univ-grenoble-alpes.fr/~demailly/manuscripts/agbook.pdf}
{Author's PDF}.
\end{thebibliography}
\end{document}
